% Compiler avec pdfLaTeX :  pdflatex lebenslauf_muster.tex  (2 fois)

\documentclass[11pt,a4paper]{article}

\usepackage[T1]{fontenc}
\usepackage[utf8]{inputenc}
\usepackage[ngerman]{babel}
\usepackage[default]{sourcesanspro} % police sans-serif professionnelle
\usepackage{microtype}
\usepackage[a4paper,left=1.8cm,right=1.8cm,top=1.2cm,bottom=1.5cm]{geometry}
\usepackage{xcolor}
\usepackage{tikz}
\usepackage{enumitem}
\usepackage[hidelinks]{hyperref}

\definecolor{shape}{HTML}{4F6FC0}
\definecolor{accent}{HTML}{5B7DB1}
\definecolor{textgray}{HTML}{3A3A3A}
\definecolor{lightgray}{HTML}{8A8A8A}

\pagestyle{empty}
\setlength{\parindent}{0pt}
\color{textgray}

% ---------- Commandes ----------
\newcommand{\sectionhead}[1]{%
  \vspace{1.1em}
  {\large\bfseries\color{accent}\MakeUppercase{#1}}\\[-0.35em]
  {\color{accent!50}\rule{\linewidth}{0.5pt}}\vspace{0.5em}}

% \entry{date}{titre}{organisation}{puces}
\newcommand{\entry}[4]{%
  \noindent
  \begin{minipage}[t]{2.9cm}\raggedright\small\color{lightgray}#1\end{minipage}%
  \hspace{0.3cm}%
  \begin{minipage}[t]{\dimexpr\linewidth-3.2cm\relax}
    {\normalsize\bfseries\color{black!90}#2}\\[-0.1em]
    {\small\color{accent}#3}
    \begin{itemize}[leftmargin=1.1em,itemsep=0.1em,topsep=0.3em,parsep=0pt,label={\tiny$\bullet$}]
      \small #4
    \end{itemize}
  \end{minipage}\par\vspace{0.55em}}

\begin{document}

% ---------- Kopfbereich ----------
\begin{tikzpicture}[remember picture,overlay]
  % Dekoratives blaues Element oben
  \fill[shape!85,rounded corners=8pt]
    ([xshift=7.0cm]current page.north west) --
    ([xshift=13.0cm]current page.north west) --
    ([xshift=11.0cm,yshift=-1.4cm]current page.north west) --
    ([xshift=8.0cm,yshift=-0.9cm]current page.north west) -- cycle;
  % Name
  \node[anchor=north west,inner sep=0pt] at ([xshift=1.8cm,yshift=-2.0cm]current page.north west)
    {\fontsize{24}{28}\selectfont\bfseries\color{black!80}Emmanuel Dakdem Nguena};
  % Kontakt
  \node[anchor=north east,inner sep=0pt,align=right,font=\footnotesize,text=black!65]
    at ([xshift=-1.8cm,yshift=-1.0cm]current page.north east)
    {Ederstraße 6, 35390 Gießen\\
     017643614045\\
     emmanuel.dakdem@yahoo.com\\
     Führerschein B197};
\end{tikzpicture}

\vspace*{2.7cm}

% Profil
{\small\linespread{1.08}\selectfont
Zielstrebiger Mechatronik-Masterstudent mit fundiertem Maschinenbau-Background und ausgeprägter
Leidenschaft für die Automatisierungstechnik. Durch die erfolgreiche Konzeption eines modularen
Odometriesystems in ROS~2 bringe ich tiefe Kenntnisse in den Bereichen Sensorfusion, Computer Vision und
autonome Systeme mit. Kombiniert mit meiner Praxiserfahrung aus der Industrie (u.\,a. Datenanalyse mit
Power~BI und SAP bei Pfeiffer Vacuum) besitze ich eine strukturierte, analytische Arbeitsweise an der
Schnittstelle zwischen Hardware und Software.\par}

\sectionhead{Berufserfahrung}

\entry{Apr. 2024 --\\ Feb. 2025}
{Werkstudent neue Technologien und Konnektivität -- Aßlar}
{PFEIFFER VACUUM}
{\item Unterstützung bei der Einführung von SAP S/4HANA
 \item Durchführung von Datenrecherchen, Scouting neuer Tools und Unterstützung bei kleineren Projekten
 \item Konzeptentwicklung zur kamerabasierten Qualitätssicherung bei einer Laserentgratanlage
 \item Strukturierte Dokumentation der Anforderungen in einem professionellen Lastenheft
 \item Abgleich der definierten Anforderungen mit bestehenden Prozessen und Systemen
 \item Koordination und Kommunikation mit potenziellen externen Lieferanten}

\entry{Apr. 2023 --\\ Apr. 2024}
{Werkstudent Qualitätskontrolle -- Heuchelheim}
{Schunk Carbon Technology GmbH}
{\item Tätig im Bereich „Qualitätssicherung“
 \item Durchführung von Messprotokollen und Dokumentation
 \item Analyse und Dokumentation der Produktionsdaten mit SAP}

\sectionhead{Ausbildung}

\entry{Okt. 2025 --\\ heute}
{Master of Science (Schwerpunkt: Mechatronik \& Robotik) -- Friedberg}
{Technische Hochschule Mittelhessen}
{\item Vertiefung in interdisziplinären Bereichen der Mechatronik und Robotik
 \item Fokus auf die Integration von Mechanik, Elektronik und Informatik; Systemsimulation und KI-Methoden}

\entry{Okt. 2021 --\\ Sep. 2025}
{Bachelor of Engineering (Schwerpunkt: Allgemeiner Maschinenbau) -- Gießen}
{Technische Hochschule Mittelhessen}
{\item Konstruktion und Inbetriebnahme technischer Maschinen und Anlagen
 \item CAD- und Simulationssoftware (Siemens NX, Ansys, SolidWorks), Automatisierungs- und Messtechnik}

\entry{Sep. 2016 --\\ Aug. 2018}
{Allgemeine Hochschulreife (Abitur) -- Yaoundé, Kamerun}
{Gymnasium CSAO}
{\item Naturwissenschaftlicher Schwerpunkt: Mathematik und Physik
 \item Fundierte Kenntnisse in analytischem Denken und technischer Problemlösung}

\sectionhead{Praktika \& Projekte}

\entry{Apr. 2026 --\\ heute}
{Projekt: Klassische und visuelle Odometrieverfahren für ein automatisiertes Modellfahrzeug}
{Technische Hochschule Mittelhessen, Friedberg}
{\item \textbf{Konzeption \& Implementierung:} Modulares Odometriesystem in ROS~2 für das Modellfahrzeug MxCarkit
 \item \textbf{Sensorfusion:} Kopplung von IMU-Daten, Motordrehzahlen und Lenkwinkel zur Zustandsschätzung und Posenberechnung
 \item \textbf{Computer Vision:} Zwei Verfahren zur visuellen Odometrie auf Basis von Stereo-Kameradaten
 \item \textbf{Systemintegration:} Einbindung in die bestehende ROS~2-Architektur über standardisierte Schnittstellen
 \item \textbf{Validierung \& Analyse:} Experimentelle Fahrversuche zur Bewertung von Genauigkeit, Robustheit und Systemgrenzen}

\entry{März 2025 --\\ Mai 2025}
{Praktikant: Analysetool zur Identifikation von Optimierungspotentialen in der Produktion}
{PFEIFFER VACUUM, Aßlar}
{\item Visualisierung von Produktionskennzahlen (z.\,B. Stückzahl, Ausschussquote) mit Power~BI
 \item Aufbereitung von Daten aus dem Dataflow, Datenmodellierung mit DAX, Erstellung von Dashboards und Berichten
 \item Analyse von Fertigungsprozessen und Identifikation von Engpässen
 \item Anwendung von SAP-Transaktionscodes (z.\,B. COOIS, CO03, MM03)}

\sectionhead{Kompetenzen}

\entry{Sprachen}{}{}
{\item Französisch (verhandlungssicher), Deutsch (verhandlungssicher), Englisch (verhandlungssicher)}

\entry{Software}{}{}
{\item Siemens NX (CAD, FEM), Python, Codesys, Siemens TIA Portal, Adams View, SAP S/4HANA, C++, ROS~2, MS Office}

\end{document}
